\guidepart{pink}{Make it yours}{Make it yours}
  {Settings, backups, and moving in from Alchementary Link.}
  {\partitem{\iGear}{Settings}
   \partitem{\iUpload}{Import and export}
   \partitem{\iPlane}{Moving house}}

% =========================================================================================
\chapter{Settings}\label{ch:settings}
\chapmeta{3 minutes}{You'll know where every setting lives}

\keys{\mod+,} opens settings (also the \iGear{} in the sidebar footer and \menu{File\then Settings}). Change things,
then click \ui{Save}.

\begin{center}\begin{minipage}[t]{0.49\linewidth}\crop{settings-general.png}{581}{253}{576}{285}{\linewidth}{General}\end{minipage}\hfill\begin{minipage}[t]{0.49\linewidth}\crop{settings-terminal.png}{581}{253}{576}{285}{\linewidth}{Terminal}\end{minipage}\end{center}

\begin{description}
  \item[General] App theme (Pro, Midnight, Light), default protocol for new profiles, whether to ask before closing or
        disconnecting, status bar on or off, sidebar width, and the local shell (leave it empty for the system default,
        or \ui{Browse…} for one).
  \item[Terminal] Colour scheme, font, size, scrollback and cursor style. These are the defaults for new profiles;
        each profile can override them.
  \item[Security] Keychain and vault status; set up, unlock or lock the vault.
  \item[Known Hosts] Saved server fingerprints, with first-seen and last-seen times.
  \item[Help] Built-in help: connection types, shortcuts and troubleshooting, when this guide isn't to hand.
  \item[About] Version, \ui{Check for Updates} and \ui{Import \textasciitilde/.ssh/config}.
\end{description}

\crop{settings-help.png}{581}{253}{576}{285}{0.8\linewidth}{The built-in help, for when this guide isn't to hand.}

\section{Updates}

Nexus Terminal checks once at launch. If there's a newer version the status bar says \ui{Update available}; click it to
see the release notes and a \ui{View release} button. Nothing downloads or installs behind your back.

% =========================================================================================
\chapter{Import and export}\label{ch:transfer}
\chapmeta{3 minutes}{You'll back up your profiles, or copy them to another computer}

\begin{tldr}
\menu{File\then Export profiles…} makes a password-protected \code{.nexus} file. \menu{File\then Import profiles…}
reads it back. Already have an \code{\textasciitilde/.ssh/config}? Import that too.
\end{tldr}

\begin{description}
  \item[Encrypted export] Choose the profiles, set a password, save a \code{.nexus} file. Saved passwords aren't
        included, so the file is safe to email to yourself.
  \item[Encrypted import] Pick a \code{.nexus} (or an old \code{.link}) file, enter its password and decide what happens to
        profiles that already exist: \ui{Skip existing}, \ui{Overwrite existing} or \ui{Create duplicates}.
  \item[JSON] \ui{Export Profiles (JSON)} and \ui{Import Profiles (JSON)} are in the palette. Tick
        \ui{Include secrets} to keep proxy passwords and TOTP secrets, which are otherwise left out.
  \item[OpenSSH config] \menu{File\then Import \textasciitilde/.ssh/config} turns your \code{Host} blocks into profiles,
        with host names, users, ports, identity files and jump hosts.
\end{description}

\crop{export.png}{894}{357}{872}{380}{0.55\linewidth}{Exporting profiles. Everything's ticked to start with; untick what you don't want.}

\begin{danger}[JSON with secrets is not encrypted]
A JSON export with \ui{Include secrets} ticked has those secrets in plain text. Use the encrypted \code{.nexus} export to
move things between machines.
\end{danger}

% =========================================================================================
\chapter{Moving house}\label{ch:move}
\chapmeta{2 minutes}{You'll know what came over from Alchementary Link and what didn't}

\begin{tldr}
It's automatic. First run copies your Alchementary Link profiles, vault and recordings. Your old folder is left
alone. Recreate trigger rules and SFTP bookmarks by hand.
\end{tldr}

When Nexus Terminal starts for the first time and finds an \code{AlchementaryLink} folder next to where its own data
goes, it \emph{copies} \code{link.db} to \code{nexus.db}, along with the \code{vault/} and \code{recordings/} folders,
and shows a toast saying where they came from.

\begin{checklist}
  \item[\color{green}\iCheck] Profiles, groups, favourites, known hosts, snippets and preferences
  \item[\color{green}\iCheck] Vault (same master password as before)
  \item[\color{green}\iCheck] Recordings
  \item[\color{green}\iCheck] Keychain passwords: picked up the first time each one is needed, then copied across
  \item[\color{green}\iCheck] Old \code{.link} export files still import
  \item[\color{rose}\ph{E4F6}] Trigger rules and SFTP bookmarks (they lived in the old app's web storage)
\end{checklist}

\begin{headsup}
If a \code{NexusTerminal/nexus.db} already exists (say, from before the Alchementary Link rebrand), it's used as it is
and nothing is copied.
\end{headsup}

\section{Where everything lives}

\renewcommand\arraystretch{1.45}
\begin{tabularx}{\linewidth}{@{}>{\leavevmode\heavy\color{ink}}p{20mm}X@{}}
  \toprule
  macOS   & \code{\textasciitilde/Library/Application Support/NexusTerminal}\\
  Windows & \code{\%APPDATA\%\textbackslash NexusTerminal}\\
  Linux   & \code{\$XDG\_CONFIG\_HOME/NexusTerminal}, usually \code{\textasciitilde/.config/NexusTerminal}\\
  \bottomrule
\end{tabularx}

Inside: \code{nexus.db} (profiles, known hosts, preferences, snippets, IRC state and the audit log),
\code{vault/vault.json}, \code{recordings/}, \code{triggers.json} and \code{sftp-bookmarks.json}. Keychain entries are
stored under the service name \code{NexusTerminal}.

\begin{protip}[Try things without touching your real data]
Start Nexus Terminal with \code{--data-dir /some/other/folder} and it uses that folder instead. Brilliant for
experiments, demos and a portable install on a USB stick.
\end{protip}
